\begin{textsample}{POS Dim 4 – am – Score 10.00 – am/am\_ef\_18.txt}  \label{ex:f4_pos_106}
5.7.
O trabalho educativo na escola perpassa por inúmeros desafios, entre eles o educar para o respeito à diversidade \textbf{religiosa}, considerando a presença do sagrado no povo brasileiro.
Um ensino voltado para a prática reflexiva e dialógica na perspectiva da construção da cidadania deve ter a preocupação e o cuidado de não fazer proselitismo e jamais permitir a prática da \textbf{intolerância}, pois os educandos, professores devem ser mutuamente reconhecidos e valorizados em sua consciência, crenças e seus credos.
O Brasil é um \textbf{Estado} \textbf{laico} e a Constituição Federal prevê a liberdade de religião, sendo a igreja e o \textbf{Estado} oficialmente separados.
A Constituição Federal de 1988 estabelece : VI - é inviolável a liberdade de consciência e de crença, sendo assegurado o livre exercício dos cultos \textbf{religiosos} e garantida, na forma da lei, a proteção aos locais de culto e a suas liturgias ; ( art. 5º, inciso VI ).
Ela ainda prevê no art. 5º, § 2º, ( “ Os direitos e garantias expressos nesta Constituição não excluem outros decorrentes do regime e dos princípios por ela adotados, ou dos tratados internacionais em que a República Federativa do Brasil seja parte ” ) da Constituição Federal do Brasil, também é aplicável o previsto no art.
XVIII da Declaração Universal dos Direitos Humanos, que expressa que : “ toda pessoa tem direito à liberdade de pensamento, consciência e religião ”, combinado com o artigo XIX, também da declaração dos direitos humanos, que expressa que “ toda pessoa tem direito à liberdade de opinião e expressão ”.
Uma educação voltada para a cidadania necessita promover o respeito às diferentes religiões e as manifestações de fé presentes na escola, pois as pessoas são livres e possuem ou não uma crença que orienta suas vidas.
A aprendizagem voltada para o respeito à alteridade é desafiante com resultados progressivos, visto que não deve ser algo imposto, mas refletido e que vise sempre à eliminação de práticas discriminatórias e antidemocráticas.
O Relatório sobre \textbf{Intolerância} e Violência \textbf{Religiosa} no Brasil ( RIVIR — 2011/2015 ), lançado em 2016, elencou oito tipos de violência a partir de relatório interno da SDH dedicado à análise dos dados da Ouvidoria Nacional dos Direitos Humanos ( SYDOW, 2015 ) : violência psicológica por motivação \textbf{religiosa} ; violência física por motivação \textbf{religiosa} ; violência relativa à prática de atos/ritos \textbf{religiosos} ; violência moral por motivação \textbf{religiosa} ; violência institucional por motivação \textbf{religiosa} ; violência patrimonial por motivação \textbf{religiosa} ; violência sexual por motivação \textbf{religiosa} e negligência por motivação \textbf{religiosa}.
O Relatório destaca ainda, que o maior peso entre os tipos de violências foi identificado entre aquelas definidas como psicológicas, seguidas da violência moral que são próximas.
Considerando o que é apontado no relatório, é importante que a educação fale sobre ética, valores humanos e cultura de paz, uma vez que historicamente a \textbf{intolerância} \textbf{religiosa} está presente desde 1500 em nosso país.
A diversidade \textbf{religiosa} no Brasil faz -nos ver que encontramos : religiões cristãs, \textbf{indígenas}, afro-brasileiras e outras.
Apesar das diferenças, nós brasileiros, podemos conviver em paz e harmonicamente visando o bem de todos.
Nesse sentido, o \textbf{chão} da escola não é \textbf{lócus} de preconceito às religiões menos valorizadas na sociedade brasileira, é sim, um espaço em que todos podem interagir e expor suas ideias e pensamentos sem preconceito e sem dogmas ou normas \textbf{religiosas}.
A UNESCO ( Organização das Nações Unidas para a Educação, Ciência e Cultura ) em sua Declaração sobre a Eliminação de Todas as formas de \textbf{Intolerância} e \textbf{Discriminação} com Base em Religião ou Crença declara que a religião e a crença devem contribuir para a paz mundial, a justiça social, a amizade entre os povos e a eliminação de ideologias ou práticas de colonialismo e \textbf{discriminação} \textbf{racial}.
A Lei Nº 11.635, de 27 de dezembro de 2007, institui o Dia Nacional de Combate à \textbf{Intolerância} \textbf{Religiosa} a ser comemorado anualmente em todo o \textbf{território} nacional no dia 21 de janeiro, a data rememora o dia do falecimento da Iyalorixá Mãe Gilda, do terreiro Axé Abassá de Ogum ( BA ), vítima de \textbf{intolerância} por ser praticante de religião de matriz africana.
Em Manaus, a Semana da Liberdade \textbf{Religiosa} é comemorada, no município de Manaus, na terceira semana de maio, conforme a Lei Nº 2146, de 05 de julho de 2016.
As escolas, os colégios e as entidades não governamentais poderão desenvolver programações, como a realização de palestras e atividades práticas de incentivo à liberdade \textbf{religiosa}.
Neste sentido, conhecer e valorizar a diversidade \textbf{religiosa} se torna como ponto essencial para se compreender e entender a necessidade de transformação da sociedade, sendo um ponto de partida no processo de ensino e de aprendizagem e jamais como ponto de chegada.
Sabemos que não é uma tarefa simples e requer um conhecimento diverso.
JUNQUEIRA et al. ( 2017 ) insere que : “ No entanto, a diversidade é o \textbf{lócus} das diferenças, ou seja, é no encontro das diferenças que a diversidade se torna reflexão entre as possíveis relações ”. ( p. 346 ), portanto, é papel da escola brasileira todo este trabalho.
O Referencial Curricular Amazonense ao discutir a temática \textbf{religiosa} se pautou em questões atuais, e os pontos a serem observados nos componentes curriculares buscam dialogar com a diversidade \textbf{religiosa} presente em nosso \textbf{Estado}.

% matched lemmas: chão, discriminação, estado, indígena, intolerância, laico, lócus, racial, religioso, território
\end{textsample}
